% \section{Alignment as a database execution pipeline}
% \label{sec:pipeline-overview}

% This is empty for now.

% BioSeqDB retains the broad minimap2 sequence of seeding, chaining, and
% extension, but places it inside a database execution pipeline. An SQL
% expression is first assigned to a PostgreSQL custom scan. The scan buffers
% compatible rows, resolves the sequence representation and target-index
% source, divides the work into memory-admitted waves, and selects a scalar
% CPU, vectorized CPU, or Vulkan route. Only then does a query enter the
% biological mapper. The mapper returns typed match records which the scan
% projects into SQL rows.

% This organization separates two related flows. The logical flow transforms
% sequence symbols into minimizers, anchors, chains, candidate regions, and
% final alignments. The physical flow decides how sequences and indexes are
% obtained, which work can share a prepared target, how much memory a wave may
% reserve, and which processor executes each supported stage. The distinction
% is important because a fast kernel does not give a fast database operator if
% the same target index is reconstructed for every row. In the same way, a GPU
% sketching kernel does not define an end-to-end GPU route if seed lookup and
% chaining still occur on the CPU.

% Figure~\ref{fig:bioseq-pipeline-dataflow} places the long-read route used by
% BioSeqDB within the conventional read-mapping flow. The upper lane shows a
% short-read pipeline from another mapper for comparison. Both lanes use the
% broad stages of seeding, chaining, extension, and traceback. BioSeqDB
% concentrates on the lower lane, where minimizers reduce the number of
% positions considered by the selected chaining route.

% \begin{figure}[tbp]
%   \centering
%   \begin{tikzpicture}[
%       x=1cm,y=0.82cm,>=Latex,font=\small,
%       panel/.style={draw=black!30,rounded corners=3pt,fill=black!6},
%       proc/.style={draw,rounded corners=2pt,fill=white,align=center,
%                    minimum height=0.66cm,text width=1.25cm,inner sep=2pt,
%                    font=\footnotesize},
%       io/.style={proc,fill=LightSteelBlue1,text width=1.15cm},
%       chip/.style={draw,rounded corners=3pt,fill=PaleGreen2,
%                    minimum width=0.52cm,minimum height=2.10cm,
%                    align=center,inner sep=1pt},
%       flow/.style={-{Latex[length=2mm]},thick}
%     ]
%     \draw[panel] (1.55,0.30) rectangle (5.10,3.62);
%     \draw[panel] (5.72,0.30) rectangle (8.70,3.62);
%     \draw[panel] (9.30,0.30) rectangle (13.08,3.62);

%     \node at (0.72,4.20) {Input};
%     \node at (3.32,4.20) {Seeding stage};
%     \node at (7.21,4.20) {Chaining stage};
%     \node at (11.19,4.20) {Extension stage};
%     \node at (14.30,4.20) {Output};

%     \node[io] (short) at (0.72,2.78) {Short reads};
%     \node[io] (long) at (0.72,1.12) {Long reads};

%     \node[proc] (smem) at (2.38,2.78) {Exact-match seeding};
%     \node[proc] (reseed) at (4.20,2.78) {Re-seeding};
%     \node[align=center,text width=1.28cm] (extract) at (2.38,1.12)
%       {Extract minimizers};
%     \node[proc] (lookup) at (4.20,1.12) {Target-index lookup};
%     \node[chip] (seeds) at (5.40,1.95) {\rotatebox{90}{Seeds}};

%     \node[proc,text width=1.55cm] (shortchain) at (6.72,2.78)
%       {Short-read\\chaining};
%     \node[proc,text width=2.00cm] (selectedchain) at (7.21,1.12)
%       {Selected chaining route};
%     \node[chip] (chains) at (9.00,1.95) {\rotatebox{90}{Chains}};

%     \node[fill=white,inner sep=1pt] at (9.78,3.60) {Generate pairs};
%     \node[proc,text width=1.55cm] (local) at (10.62,2.78)
%       {Local\\extension};
%     \node[proc,text width=2.05cm] (full) at (10.62,1.12)
%       {Optional base alignment};
%     \node[draw,rounded corners=2pt,fill=white,
%           minimum width=0.66cm,minimum height=2.10cm,align=center]
%       (trace) at (12.72,1.95) {\rotatebox{90}{Traceback}};

%     \node[io,text width=1.50cm] (output) at (14.30,1.95)
%       {Positions\\edit paths};

%     \draw[flow] (short) -- (smem);
%     \draw[flow] (smem) -- (reseed);
%     \draw[flow] (long) -- (extract);
%     \draw[flow] (extract) -- (lookup);
%     \draw[flow] (reseed.east) -- ($(seeds.west)+(0,0.68cm)$);
%     \draw[flow] (lookup.east) -- ($(seeds.west)+(0,-0.68cm)$);
%     \draw[flow] ($(seeds.east)+(0,0.68cm)$) -- (shortchain.west);
%     \draw[flow] ($(seeds.east)+(0,-0.68cm)$) -- (selectedchain.west);
%     \draw[flow] (shortchain.east) -- ($(chains.west)+(0,0.68cm)$);
%     \draw[flow] (selectedchain.east) -- ($(chains.west)+(0,-0.68cm)$);
%     \draw[flow] ($(chains.east)+(0,0.68cm)$) -- (local.west);
%     \draw[flow] ($(chains.east)+(0,-0.68cm)$) -- (full.west);
%     \draw[flow] (local.east) -- ($(trace.west)+(0,0.68cm)$);
%     \draw[<->,thick] (full.east) --
%       node[below,yshift=-1pt] {z-drop} ($(trace.west)+(0,-0.68cm)$);
%     \draw[flow] (trace) -- (output);
%   \end{tikzpicture}
%   \caption{Read-alignment stages from input to reported coordinates and edit
%   paths. The upper lane is an external short-read comparison. The lower lane
%   is BioSeqDB's minimizer-based route. The lanes follow the same broad stages
%   without sharing literal interfaces.}
%   \label{fig:bioseq-pipeline-dataflow}
% \end{figure}

% \section{The logical alignment chain}
% \label{sec:pipeline-logical-chain}

% The generic CPU and SIMD mapper is assembled as a typed push pipeline. Each
% stage accepts a batch view and emits the next representation while retaining
% any required data or index owner. The stage type also fixes the selected
% algorithm. Runtime choices are resolved before the main loop, after which each
% batch follows one concrete route. The forwarding contract requires exactly
% one downstream emission, which prevents a stage from dropping or duplicating
% a batch.

% This pipeline has two entry points. The usual entry begins with sequence
% symbols and computes query minimizers. A second entry accepts authenticated
% prebuilt minimizers and starts at the index hand-off. The second form is used
% only when the minimizer representation, source identity, alphabet, \(w\), and
% \(k\) match the requested operation. It is a structural omission of the
% sketch stage, not an empty sketch operation.

% Algorithm~\ref{alg:bioseq-logical-pipeline} gives its common order. Optional
% work is governed by the SQL projection and predicates. Divergence and mapping
% quality are not calculated when neither projection nor filtering needs them.
% Base-level extension is skipped when chain coordinates and scores are
% sufficient.

% \begin{algorithm}[tbp]
%   \caption{Logical mapping of one query against a prepared target}
%   \label{alg:bioseq-logical-pipeline}
%   \small
%   \begin{algorithmic}[1]
%     \Require Query \(Q\), target \(T\), target index \(I\), parameters \(P\)
%     \Ensure Ordered mapping results
%     \State \textbf{if} the entry point begins with sequence symbols \textbf{then}
%     \State \hspace{\algorithmicindent}
%       \(M \gets \operatorname{Sketch}(Q,P.w,P.k,P.\mathit{minimizer})\)
%     \State \textbf{else}
%     \State \hspace{\algorithmicindent}
%       \(M \gets \operatorname{ValidatedPrebuiltMinimizers}(Q,P)\)
%     \State \textbf{end if}
%     \State \(A \gets \operatorname{LookupAndCanonicalize}(M,I,P.\mathit{maxOcc})\)
%     \State \((C,A') \gets \operatorname{Chain}(A,P.\mathit{chainer})\)
%     \State \(R \gets \operatorname{GenerateRegions}(C,A',|Q|)\)
%     \State \(R \gets \operatorname{ClassifyAndReduce}(R,P)\)
%     \State \textbf{if} divergence is required \textbf{then}
%     \State \hspace{\algorithmicindent}
%       \(R \gets \operatorname{EstimateDivergence}(R,M,A',P.k)\)
%     \State \textbf{end if}
%     \State \textbf{if} an exact edit predicate permits conservative prefiltering
%       \textbf{then}
%     \State \hspace{\algorithmicindent}
%       \(R \gets \operatorname{RejectImpossibleGapJobs}(R,A',Q,T,P)\)
%     \State \textbf{end if}
%     \State \textbf{if} base-level fields or an exact edit predicate require
%       extension \textbf{then}
%     \State \hspace{\algorithmicindent}
%       \(R \gets \operatorname{Extend}(R,A',Q,T,P.\mathit{aligner})\)
%     \State \textbf{end if}
%     \State \textbf{if} an exact edit predicate is active \textbf{then}
%     \State \hspace{\algorithmicindent}
%       \(R \gets \operatorname{ApplyExactEditPredicate}(R,P)\)
%     \State \textbf{end if}
%     \State \(R \gets \operatorname{ClassifyAndReduce}(R,P)\)
%     \State \textbf{if} mapping quality is required \textbf{then}
%     \State \hspace{\algorithmicindent}
%       \(R \gets \operatorname{ComputeMapq}(R,P)\)
%     \State \textbf{end if}
%     \State \Return \(\operatorname{ToMatchResults}(R,|Q|,|T|)\)
%   \end{algorithmic}
% \end{algorithm}

% \subsection{Target preparation and minimizers}

% Target preparation occurs outside the per-query chain. A target may be an
% in-memory compact sequence, a lazy PostgreSQL sequence, a fragmented value,
% or a persisted minimizer-index generation. The scheduler first resolves
% whether the index will be recomputed or read from persistent storage. A
% prepared CPU target or resident GPU target can then be shared by all
% compatible queries in a dispatch wave. This placement removes target
% sketching and index construction from repeated query work.

% A persisted index is accepted only when its generation, source identity,
% storage geometry, alphabet, and minimizer parameters match the request. The
% pipeline may decode bounded fragments or select partitions, but seed lookup
% must observe complete occurrence counts and position lists for every
% minimizer it consumes. A fragment-native CPU or SIMD preparation can go one
% step further and produce ordered anchors for each query. The later seed stage
% then performs a checked zero-copy hand-off while the owning batch remains
% alive.

% When minimizers are not prebuilt, the minimizer stage reads the lazy query
% view and invokes the scalar, SIMD, or Vulkan extractor selected by its trait.
% CPU implementations write into the current arena. The Vulkan form can return
% a host view together with a device-resident span, which lets GPU seed lookup
% reuse the sketch without another upload. An empty query or a query shorter
% than \(k\) produces a valid empty minimizer span.

% \subsection{Seeds, chains, and candidate regions}

% Seed lookup maps each query minimizer to target occurrences. It discards
% absent and over-repetitive values according to the occurrence threshold,
% rejects incomplete persisted postings, converts reverse-strand positions to
% a common query coordinate, and establishes one canonical anchor order. The
% position domain remains 64 bit so that a large reference is not narrowed at
% the stage boundary. Every chaining implementation receives this same order.

% The chain stage selects scalar or SIMD dynamic programming, range-minimum-query
% chaining, or tiled Vulkan dynamic programming. It returns
% chains and a span of anchors reordered into chain order. Region generation
% then produces one candidate region for each retained chain, sorts candidates
% by score, and derives half-open reference and strand-correct query
% coordinates from the first and last anchors. It also records seeded match and
% block lengths used by later heuristics.

% Figure~\ref{fig:pipeline-seed-chain-extend} shows the transition on a small
% example. Exact two-base seeds locate several occurrences in the reference.
% Chaining compares their order and displacement, then retains the collinear
% pair. Dynamic programming fills the sequence interval between the anchors.
% The example uses a compact score matrix to make the extension boundary and
% traceback explicit.

% \begin{figure}[tbp]
%   \centering
%   \begin{tikzpicture}[
%       x=0.69cm,y=0.72cm,>=Latex,font=\small,
%       cell/.style={draw,minimum width=0.39cm,minimum height=0.39cm,
%                    inner sep=0pt,align=center},
%       tile/.style={draw,minimum width=0.65cm,minimum height=0.42cm,
%                    inner sep=1pt,align=center},
%       seedA/.style={cell,fill=PaleGreen2},
%       seedB/.style={cell,fill=Thistle2},
%       tileA/.style={tile,fill=PaleGreen2},
%       tileB/.style={tile,fill=Thistle2},
%       neutral/.style={cell,fill=white},
%       extended/.style={cell,draw=DeepSkyBlue2,very thick,fill=white},
%       occA/.style={draw=ForestGreen,circle,fill=PaleGreen2,
%                    minimum size=0.34cm,inner sep=0pt},
%       occB/.style={draw=Orchid3,circle,fill=Thistle2,
%                    minimum size=0.34cm,inner sep=0pt},
%       dp/.style={draw,minimum width=0.56cm,minimum height=0.52cm,
%                  inner sep=0pt,align=center},
%       flow/.style={-{Latex[length=1.8mm]},thick},
%       distarrow/.style={-{Latex[length=1.8mm]},thick,draw=DeepSkyBlue3,
%                         text=DeepSkyBlue4},
%       trace/.style={-{Latex[length=1.6mm]},very thick,draw=ForestGreen},
%       panelno/.style={draw,circle,fill=black,text=white,
%                       minimum size=0.42cm,inner sep=0pt}
%     ]
%     \draw[dashed,black!45] (6.80,0.05) -- (6.80,9.25);
%     \draw[dashed,black!45] (13.40,0.05) -- (13.40,9.25);

%     \node at (2.62,8.98) {Query};
%     \foreach \x/\base/\style in
%       {0/T/seedA,1/A/seedA,2/A/neutral,3/G/neutral,
%        4/A/neutral,5/T/neutral,6/C/seedB,7/T/seedB} {
%       \node[\style] at ({0.70+0.55*\x},8.32) {\base};
%     }
%     \draw[distarrow] (0.98,7.62) --
%       node[below] {\(d=4\)} (3.18,7.62);
%     \node[align=center] at (2.32,7.02) {Identify reference positions};
%     \draw[flow] (2.32,6.72) -- (2.32,6.42);

%     \draw[rounded corners=3pt] (0.30,3.40) rectangle (4.25,6.40);
%     \node at (2.28,6.03) {Reference index};
%     \node[tileA] at (1.05,5.35) {TA};
%     \node[tile] at (2.28,5.35) {GC};
%     \node[tileB] at (3.51,5.35) {CT};
%     \node[occA] at (0.87,4.55) {0};
%     \node[occA] at (1.23,4.10) {8};
%     \node at (1.05,3.66) {\(\vdots\)};
%     \node at (2.28,4.55) {2};
%     \node at (2.28,4.10) {5};
%     \node at (2.28,3.66) {\(\vdots\)};
%     \node[occB] at (3.33,4.55) {3};
%     \node[occB] at (3.69,4.10) {6};
%     \node at (3.51,3.66) {\(\vdots\)};
%     \draw[flow] (4.25,4.92) -- (5.13,4.92);

%     \node[rotate=90] at (4.85,3.93) {Reference};
%     \foreach \y/\base/\style in
%       {0/T/seedA,1/A/seedA,2/G/neutral,3/C/seedB,4/T/seedB,
%        5/G/neutral,6/C/seedB,7/T/seedB,8/T/seedA,9/A/seedA} {
%       \node[\style] at (5.45,{6.18-0.49*\y}) {\base};
%       \node[anchor=west] at (5.78,{6.18-0.49*\y}) {\y};
%     }
%     \node[panelno] at (2.13,0.38) {1};
%     \node[anchor=west] at (2.50,0.38) {Seed};

%     \node at (10.10,8.65) {Chain candidates};
%     \node at (7.55,7.82) {(a)};
%     \node[tileA] (a1) at (7.55,6.95) {TA};
%     \node[tileB] (a2) at (7.55,4.75) {CT};
%     \draw[distarrow] (a1) -- (a2);

%     \node at (8.95,7.82) {(b)};
%     \node[tileB] (b1) at (8.95,6.35) {CT};
%     \node[tileA] (b2) at (8.95,4.10) {TA};
%     \draw[distarrow] (b2) -- (b1);

%     \node at (10.35,7.82) {(c)};
%     \node[tileB] at (10.35,5.98) {CT};
%     \node[tileA] at (10.35,5.10) {TA};

%     \draw[draw=none,fill=Burlywood1,rounded corners=4pt]
%       (11.65,3.35) rectangle (12.78,7.60);
%     \node at (12.22,7.82) {(d)};
%     \node[tileA] (d1) at (12.22,6.95) {TA};
%     \node[tileB] (d2) at (12.22,4.75) {CT};
%     \draw[distarrow] (d1) -- node[left] {\(d'=4\)} (d2);
%     \draw[thick] (12.68,6.95) -- (13.25,6.95);
%     \draw[thick] (12.68,4.75) -- (13.25,4.75);
%     \node[panelno] at (9.65,0.38) {2};
%     \node[anchor=west] at (10.02,0.38) {Chain};

%     \node at (16.18,8.98) {Query};
%     \foreach \x/\base/\style in
%       {0/T/seedA,1/A/seedA,2/A/extended,3/G/extended,
%        4/A/extended,5/T/extended,6/C/seedB,7/T/seedB} {
%       \node[\style] at ({14.25+0.55*\x},8.32) {\base};
%     }

%     \node[rotate=90] at (13.76,5.18) {Reference};
%     \foreach \y/\base/\style in
%       {0/T/seedA,1/A/seedA,2/G/extended,3/C/extended,4/T/extended,
%        5/G/extended,6/C/seedB,7/T/seedB,8/T/seedA,9/A/seedA} {
%       \node[\style] at (14.32,{7.25-0.49*\y}) {\base};
%     }

%     \node[text=DeepSkyBlue4] at (17.69,7.70) {Extended interval};
%     \foreach \c/\base in {1/A,2/G,3/A,4/T} {
%       \node at ({16.05+0.82*\c},7.15) {\base};
%     }
%     \foreach \r/\base in {1/G,2/C,3/T,4/G} {
%       \node at (15.45,{6.62-0.82*\r}) {\base};
%     }
%     \foreach \r/\c/\score/\fillname in {
%       0/0/-6/LightSteelBlue1,0/1/-9/LightSteelBlue1,
%       0/2/-12/black!15,0/3/-15/black!15,0/4/-18/black!15,
%       1/0/-9/black!15,1/1/-9/white,1/2/-7/LightSteelBlue1,
%       1/3/-10/white,1/4/-13/white,
%       2/0/-12/black!15,2/1/-12/white,2/2/-10/white,
%       2/3/-10/LightSteelBlue1,2/4/-13/white,
%       3/0/-15/black!15,3/1/-15/white,3/2/-13/white,
%       3/3/-13/white,3/4/-8/LightSteelBlue1,
%       4/0/-18/black!15,4/1/-18/white,4/2/-13/white,
%       4/3/-16/white,4/4/-11/LightSteelBlue1} {
%       \node[dp,fill=\fillname] (m\r\c)
%         at ({16.05+0.82*\c},{6.62-0.82*\r}) {\(\score\)};
%     }
%     \draw[DeepSkyBlue3,very thick]
%       (m11.north west) rectangle (m44.south east);
%     \draw[trace] (m44.north) -- (m34.south);
%     \draw[trace] (m34.north west) -- (m23.south east);
%     \draw[trace] (m23.north west) -- (m12.south east);
%     \draw[trace] (m12.north west) -- (m01.south east);
%     \draw[trace] (m01.west) -- (m00.east);

%     \draw[DeepSkyBlue3,very thick] (15.72,2.14) rectangle (16.12,2.54);
%     \node[anchor=west,font=\footnotesize] at (16.25,2.34) {extended interval};
%     \draw[trace] (15.72,1.73) -- (16.14,1.73);
%     \node[anchor=west,font=\footnotesize] at (16.25,1.73) {traceback};
%     \node[align=center,text width=4.70cm,font=\footnotesize] at (18.00,1.15)
%       {scores: \(+2\) match, \(-3\) mismatch or gap};
%     \node[panelno] at (17.15,0.38) {3};
%     \node[anchor=west] at (17.52,0.38) {Extend};
%   \end{tikzpicture}
%   \caption{Worked seed, chain, and extend example. Two exact seeds identify
%   candidate reference positions. The selected pair has consistent order and
%   displacement. Dynamic programming then aligns the interval
%   between the anchors, with the green arrows marking traceback.}
%   \label{fig:pipeline-seed-chain-extend}
% \end{figure}

% \subsection{Reduction, extension, and output}

% Candidate reduction first assigns primary and secondary relationships among
% overlapping regions, then compacts the survivor prefix according to score and
% overlap parameters. Divergence estimation is a demand-driven annotation
% derived from seed coverage. Neither operation changes the coordinate
% conventions established by region generation.

% The optional prefilter is a BioSeqDB addition. For a region with an active
% exact edit predicate, it constructs mandatory seed-to-seed gap jobs and uses
% SneakySnake or GateKeeper to obtain a conservative lower bound on the number
% of edits. A region is rejected only when the accumulated lower bound already
% exceeds the predicate threshold. Unsupported geometry proceeds to exact
% alignment. The prefilter therefore avoids some extension work without
% replacing the chosen-path edit count.

% The extension stage selects KSW2, Block Aligner, an inter-region batched KSW2
% route, or an admitted Vulkan CIGAR route. CPU batching groups short fill
% problems by length, bandwidth, and orientation, while long fills use an
% intra-sequence kernel. The Vulkan route forms device jobs and chooses a short
% script or bounded dynamic programming with traceback according to the admitted
% shape. Packed edit operations are decoded into owned result fields after
% device execution.

% After extension, any exact edit predicate is applied. The primary and
% secondary classifier then runs again. If filtering removed a candidate,
% reclassification ensures that parent choice, secondary scores, and mapping
% quality refer only to the final candidate set.
% The sink converts each region to a match record with 64-bit coordinates,
% scores, seed and chain metadata, and only the optional divergence, mapping
% quality, edit count, or CIGAR fields required by projection or filtering.

% Table~\ref{tab:pipeline-stage-contracts-a} summarizes the stage boundaries.
% The table groups adjacent control stages where they share one semantic
% contract.

% \begin{table}[tbp]
%   \centering
%   \small
%   \caption{Contracts of the BioSeqDB logical alignment pipeline.}
%   \label{tab:pipeline-stage-contracts-a}
%   \setlength{\tabcolsep}{4pt}
%   \begin{tabularx}{\linewidth}{
%       >{\raggedright\arraybackslash}p{0.14\linewidth}
%       >{\raggedright\arraybackslash}p{0.18\linewidth}
%       >{\raggedright\arraybackslash}p{0.21\linewidth}
%       >{\raggedright\arraybackslash}X}
%     \toprule
%     Stage & Input & Output & Required property \\
%     \midrule
%     Minimizer & lazy query symbols & host minimizers and optional device span &
%     Deterministic representation for \(w\), \(k\), alphabet, and position
%     domain. Short queries yield an empty valid span. \\
%     Index hand-off & query minimizers & minimizers with a target-index reference &
%     The index identity and lifetime match the target and cannot alias another
%     segment. \\
%     Seed lookup & minimizers and index & canonical anchors &
%     Occurrence pruning uses complete postings and preserves strand-correct
%     full-width coordinates. \\
%     Chaining & canonical anchors & chains and reordered anchors &
%     Every kernel consumes the same scoring parameters and anchor order. \\
%     Region reduction & chains and anchors & classified candidate prefix &
%     Coordinates are valid and removed candidates are not visible downstream. \\
%     Prefilter & candidates and sequence windows & conservative survivor prefix &
%     Only a lower bound above the exact threshold can reject a candidate. \\
%     Extension and exact filter & candidates, anchors, and sequences & scored
%     regions with optional edit paths &
%     The selected aligner supplies alignment-derived fields, including the
%     chosen-path edit count. \\
%     Reclassification and sink & post-filter candidates & ordered match records &
%     Final classification excludes removed regions and optional fields are
%     computed only when required. \\
%     \bottomrule
%   \end{tabularx}
% \end{table}

% \section{Compile-time algorithm selection and runtime routing}
% \label{sec:pipeline-dispatch}

% The runtime parameter object contains minimizer dimensions, scoring limits,
% output requirements, prefilter policy, and algorithm selectors. The
% dispatcher resolves these values into a tuple containing the sequence shape,
% base aligner, chainer, minimizer implementation, and pipeline entry point.
% Registered common combinations receive compile-time specializations. Other
% valid combinations use a typed generic implementation. This arrangement
% removes repeated selector branches from high-frequency stage calls without
% instantiating every possible product of \(w\), \(k\), sequence length, and
% algorithm choice.

% Scalar and SIMD execution share the same logical pipeline. A SIMD selection
% replaces the relevant minimizer, chaining, or alignment kernel, while region
% generation and candidate classification remain scalar control stages. An
% ``all SIMD'' route therefore means that every biological stage with a SIMD
% implementation used it. It does not imply that input handling, allocation,
% or result construction consists of vector instructions.

% CPU concurrency is controlled at statement scope. Independent queries can be
% distributed through a oneAPI Threading Building Blocks (TBB) task arena.
% Coarse inner work, such as alignment
% regions or fragment preparation, may also be parallelized, but only one
% nested coarse region is active at a time. A nested request that would create
% another fan-out runs serially. This rule bounds the effective worker count and
% prevents a query-level parallel loop from multiplying a region-level loop.

% The dedicated Vulkan route has a narrower contract. A full GPU selection
% requires compatible Vulkan minimizer extraction, resident index lookup,
% tiled chaining, and coordinate reduction. CIGAR output additionally requires
% an admitted device extension shape. A GPU work wave contains at least two
% jobs so that fixed submission costs are not charged to a single comparison.
% If the runtime, target representation, or requested output does not satisfy
% the selected mandatory route, execution returns an error. It does not perform
% the biological operation on a CPU while retaining the GPU selection.

% A standalone Vulkan prefilter does not select the full GPU route. It executes
% one conservative filter inside an otherwise CPU or SIMD pipeline. This
% distinction prevents a device-accelerated filter from changing the meaning of
% the mapper's backend label.

% Persisted GPU indexes also have a specific boundary. A compatible resident
% partition can support direct coordinate projection without rebuilding target
% minimizers. Persisted CIGAR projection is not supported because extension
% needs target windows in a device-readable representation. The pipeline
% rejects that combination instead of hydrating the target and completing the
% alignment on the CPU.

% \begin{table}[tbp]
%   \centering
%   \small
%   \caption{Selectable biological stages and their execution routes.}
%   \label{tab:pipeline-algorithm-axes}
%   \begin{tabularx}{\linewidth}{
%       >{\raggedright\arraybackslash}p{0.21\linewidth}
%       >{\raggedright\arraybackslash}X}
%     \toprule
%     Stage & Available routes \\
%     \midrule
%     Minimizer extraction & scalar minimap-compatible sketching, SIMD
%     rolling-hash variants, and Vulkan minimizer extraction \\
%     Seed lookup & in-memory or persisted CPU providers, fragment-native
%     prepared anchors, and resident Vulkan index lookup \\
%     Chaining & scalar dynamic programming, SIMD dynamic programming,
%     range-minimum-query chaining, and tiled Vulkan dynamic programming \\
%     Final alignment & KSW2, Block Aligner, inter-region batched KSW2, and an
%     admitted Vulkan CIGAR route \\
%     Prefilter & scalar or SIMD SneakySnake, Vulkan SneakySnake, and Vulkan
%     GateKeeper within their supported input domains \\
%     \bottomrule
%   \end{tabularx}
% \end{table}

% \section{Batching, memory, and target reuse}
% \label{sec:pipeline-batching}

% The PostgreSQL custom scan does not invoke the mapper once for each input
% tuple. Its dispatch buffer retains eligible scan slots together with their
% query and target operands. The buffer is drained when it reaches its bounded
% capacity or when input ends. Each batch is homogeneous in alphabet and packed
% storage width, so the typed processor resolves that pair once instead of
% switching representation in the per-row loop. During the read phase, result
% cursors restore the association between a mapping result and its buffered SQL
% row.

% Rows are grouped by query mode, physical target identity, and index policy.
% This grouping also supports a broadcast query, a query expression evaluated
% per row, and several independent query predicates. When several predicates
% use the same active wave, their matching row identifiers are intersected.
% The physical work can therefore share target preparation without changing the
% logical conjunction expressed by SQL.

% Sequence operands may be flat, fragmented, or lazy. Some kernels consume a
% lazy compact range directly, while others require a flat payload. Before a
% wave is dispatched, the scheduler estimates its required hydration bytes and
% requests a memory permit. If the reservation succeeds, only operands required
% by the selected kernel are materialized. Their storage is destroyed before
% the permit is released. If hydration is not admitted and the kernel supports
% streaming, the operands remain lazy and are read through the PostgreSQL
% scanner. Reducing materialization does not change the selected biological
% algorithm.

% The memory controller combines the process footprint, current system
% availability, a process soft limit, a minimum free-memory reserve, and any
% index-specific limit. Parallel work first attempts the requested fan-out,
% then reduces the wave size until it fits. If one parallel unit cannot be
% admitted, the same algorithm runs serially when its minimum working set fits.
% This is a scheduling fallback, not a scalar algorithm substitution.

% Prepared targets define the main reuse scope. A CPU target index, persisted
% reader, fragment-native anchor owner, or resident Vulkan index is constructed
% once for a compatible group of queries. Query-local arena scopes are released
% after each serial query, while the outer prepared target remains alive for
% the wave. On a CPU, sufficiently large groups of shareable queries may run in
% parallel. A selected Vulkan route suppresses this CPU fan-out because the
% device scheduler must observe the complete work set and form resident jobs.
% Both routes return results in input-query order.

% Algorithm~\ref{alg:bioseq-batch-dispatch} states the wave-level policy. The
% essential ordering is admission before materialization, target preparation
% before query execution, and result scattering before wave storage is
% released.

% \begin{algorithm}[tbp]
%   \caption{Memory-aware dispatch of buffered SQL rows}
%   \label{alg:bioseq-batch-dispatch}
%   \small
%   \begin{algorithmic}[1]
%     \Require Buffered rows \(B\), typed sequence layout \(S\), parameters \(P\)
%     \Ensure Mapping results associated with their input row identifiers
%     \State Group \(B\) by query mode, target identity, and index policy
%     \State Estimate materialization bytes for each group
%     \State Partition groups into waves under the live memory limit
%     \State \textbf{for each} wave \(W\) \textbf{do}
%       \State \(permit \gets \operatorname{TryReserve}(\operatorname{Bytes}(W))\)
%       \State \textbf{if} \(permit\) exists \textbf{then}
%       \State \hspace{\algorithmicindent}Materialize only operands required by
%         the selected route
%       \State \textbf{else}
%       \State \hspace{\algorithmicindent}Retain lazy operands if the selected
%         route supports streaming
%       \State \textbf{end if}
%       \State Prepare or reuse one target object for each target identity
%       \State \textbf{if} a mandatory Vulkan tuple is selected \textbf{then}
%       \State \hspace{\algorithmicindent}Validate wave size, runtime support,
%         index form, and output shape
%       \State \hspace{\algorithmicindent}Run resident query jobs on the device
%       \State \textbf{else}
%       \State \hspace{\algorithmicindent}Run typed scalar or SIMD queries
%         serially or in the TBB arena
%       \State \textbf{end if}
%       \State Scatter result collections to buffered row identifiers
%       \State Destroy wave storage and release \(permit\)
%     \State \textbf{end for}
%   \end{algorithmic}
% \end{algorithm}

% \section{Relation to minimap2 and GPU mappers}
% \label{sec:pipeline-comparison}

% Upstream minimap2~\cite{minimap2_paper} organizes long-read mapping as an
% imperative sequence. It sketches the query, retrieves minimizer occurrences
% from a reference index, prepares and sorts anchors, applies dynamic-programming
% or range-minimum-query chaining, generates regions, performs post-chain
% processing and extension, and calculates mapping quality. Parallelism is
% mainly applied across input fragments, with a worker-local allocation buffer.
% Figure~\ref{fig:pipeline-minimap-workflow} presents this three-module structure
% and the main substitution points used by vectorized variants.

% \begin{figure}[tbp]
%   \centering
%   \begin{tikzpicture}[
%       x=0.80cm,y=0.78cm,>=Latex,font=\small,
%       io/.style={draw,trapezium,trapezium left angle=70,
%                  trapezium right angle=110,fill=white,align=center,
%                  minimum height=0.72cm,text width=1.30cm,inner sep=2pt},
%       module/.style={draw,fill=white,align=center,minimum height=1.15cm,
%                      text width=1.52cm,inner sep=2pt},
%       inner/.style={draw,rounded corners=2pt,fill=white,align=center,
%                     minimum height=0.70cm,text width=1.22cm,inner sep=2pt},
%       alt/.style={draw,rounded corners=2pt,fill=white,align=center,
%                   minimum height=0.84cm,text width=1.92cm,inner sep=2pt},
%       flow/.style={-{Latex[length=1.9mm]},thick}
%     ]
%     \draw[Blue,thin] (-0.35,0.10) rectangle (18.45,8.15);

%     \node[io] (index) at (1.30,5.95) {Build or load index};
%     \node[io] (reads) at (3.45,5.95) {Load read sequences};

%     \draw (4.55,4.80) rectangle (10.42,7.45);
%     \node[anchor=west] at (4.75,7.08) {\textbf{Seeding}};
%     \node[inner] (extractm) at (5.45,5.82) {Extract minimizers};
%     \node[inner] (collectm) at (7.46,5.82) {Collect hits};
%     \node[inner] (sorta) at (9.45,5.82) {Prepare and sort anchors};

%     \node[module] (chain) at (11.82,5.82)
%       {\textbf{Chaining}\\dynamic programming};
%     \node[module] (align) at (14.35,5.82)
%       {\textbf{Alignment}\\SSE2 base alignment};
%     \node[io] (mapped) at (17.10,5.82) {Mapping output};

%     \draw[flow] (index) -- (reads);
%     \draw[flow] (reads) -- (extractm);
%     \draw[flow] (extractm) -- (collectm);
%     \draw[flow] (collectm) -- (sorta);
%     \draw[flow] (sorta) -- (chain);
%     \draw[flow] (chain) -- (align);
%     \draw[flow] (align) -- (mapped);

%     \node[alt] (hash) at (4.20,3.28) {Hash-table minimizer lookup};
%     \draw[<->,thick] (hash.north east) -- (collectm.south west);

%     \draw[Blue,dashed,thick,rounded corners=2pt]
%       (6.15,1.35) rectangle (15.85,4.15);
%     \node[Blue,fill=white,inner sep=1.5pt] at (10.75,3.82)
%       {Optimization alternatives};
%     \node[alt] (learned) at (7.55,2.45)
%       {Learned-index minimizer lookup};
%     \node[alt] (vchain) at (11.45,2.45)
%       {Vectorized chaining with AVX-512};
%     \node[alt] (valign) at (14.48,2.45)
%       {AVX-512 base alignment};

%     \draw[<->,very thick] (learned.north) -- (collectm.south);
%     \draw[flow] (vchain.north) -- (chain.south);
%     \draw[flow] (valign.north) -- (align.south);
%     \node[align=center,text width=14.0cm] at (9.25,0.63)
%       {Minimap2 separates seeding, chaining, and alignment.
%        Alternative lookup and vectorized kernels attach at these stage
%        boundaries.};
%   \end{tikzpicture}
%   \caption{The minimap2 mapping workflow and its principal replacement
%   points. The main line contains seeding, chaining, and base-level alignment.
%   The lower boxes show alternative lookup and vectorized kernels attached to
%   the same interfaces. They illustrate modular optimization choices and do
%   not denote BioSeqDB-specific implementations.}
%   \label{fig:pipeline-minimap-workflow}
% \end{figure}

% BioSeqDB preserves this recognizable algorithmic backbone but changes the
% execution context. Queries are packed, alphabet-aware ranges rather than only
% text strings. The target index may be in memory, persisted with a database
% generation, selectively decoded, or resident on a device. Runtime selectors
% are resolved into typed stages before query execution. Work is grouped by a
% custom scan rather than by a command-line reader. Parallelism is controlled
% by both a TBB arena and memory admission. Results are typed records inserted
% into a relational projection rather than only mapping-format lines.

% The scope is also narrower in some biological dimensions. BioSeqDB focuses on
% database-resident long-read and single-target mapping. It does not claim all
% paired-read, splice, rescue, or multi-segment modes supported by minimap2.
% The pipeline adds database-oriented features instead, including lazy sequence
% views, persistent index identity, projection-driven optional fields,
% conservative edit prefiltering, and reuse of one prepared target across many
% SQL rows.

% G3SA~\cite{g3sa_paper} provides a useful comparison for accelerator design.
% Its minimap2 route treats sketching, seed lookup, chaining, and extension as a
% resident GPU work graph. Anchors are stored in a device-friendly layout,
% tiled chaining reduces repeated global-memory access, and gridded traceback
% stores boundaries instead of a complete direction matrix. These choices show
% why batching and residency matter at least as much as an isolated fast
% kernel.

% BioSeqDB applies the same systems observation through Vulkan, but database
% execution imposes different constraints. Query batches are grouped by exact
% target identity. A target index can be retained across jobs and intermediate
% device spans can pass directly between sketching and seed lookup. Tiled
% chaining produces candidate coordinates, while an admitted extension route
% can generate packed CIGAR operations. Host code still constructs the final
% SQL-visible objects. Memory permits bound host materialization and device work
% is divided into bounded waves.

% The failure policy also differs. A GPU mapper may send exceptional reads to a
% CPU worker. BioSeqDB does not silently do so after a mandatory Vulkan route
% has been selected for an SQL statement. Unsupported tuples, one-job waves,
% unavailable device operations, and incompatible persisted output shapes are
% reported as errors. The scheduler may choose a CPU or SIMD route before it
% enters the mandatory device path. Once the device route is selected, backend
% substitution would change the stated execution contract.

% \begin{table}[tbp]
%   \centering
%   \small
%   \caption{Pipeline-level differences between minimap2, G3SA, and BioSeqDB.}
%   \label{tab:pipeline-system-comparison}
%   \setlength{\tabcolsep}{4pt}
%   \begin{tabularx}{\linewidth}{
%       >{\raggedright\arraybackslash}p{0.15\linewidth}
%       >{\raggedright\arraybackslash}p{0.23\linewidth}
%       >{\raggedright\arraybackslash}p{0.23\linewidth}
%       >{\raggedright\arraybackslash}X}
%     \toprule
%     Aspect & minimap2 & G3SA minimap2 route & BioSeqDB \\
%     \midrule
%     Work unit & input read fragment & fixed GPU read batch & PostgreSQL
%     custom-scan wave \\
%     Sequence form & text query with an external reference index & pinned host
%     input and device encodings & packed lazy, fragmented, or scanner-backed
%     ranges \\
%     Index lifetime & shared by worker calls & retained on the GPU & in-memory,
%     persisted, fragment-native, or resident with database identity \\
%     Control flow & imperative mapper with runtime options & stage-specific GPU
%     batches & runtime-to-template routing followed by typed push stages \\
%     Parallelism & worker loop across fragments & GPU batches with host support
%     threads & TBB tasks and Vulkan waves constrained by memory admission \\
%     GPU failure & no upstream GPU route & exceptional work may move to CPU & a
%     selected mandatory Vulkan route fails if it cannot execute \\
%     Result surface & mapping-format records & packed alignments and mapping
%     output & typed match fields projected into SQL rows \\
%     \bottomrule
%   \end{tabularx}
% \end{table}

% \section{Final invariants and current boundaries}
% \label{sec:pipeline-final}

% The resulting pipeline is governed by six main invariants.

% \begin{enumerate}
%   \item The generic CPU and SIMD mapper uses one typed stage order. The
%   resident GPU route uses a separate graph.
%   \item Runtime algorithm selectors are consumed before the hot loop. Common
%   tuples are specialized and valid uncommon tuples use a typed fallback.
%   \item Arena spans, lazy providers, persisted readers, and device handles
%   carry ownership that outlives downstream use.
%   \item Prepared targets have an explicit reuse scope across compatible query
%   waves. Query-local storage can be released without invalidating the target.
%   \item Worker concurrency and memory admission jointly bound parallel work.
%   Reducing fan-out or materialization does not change the biological kernel.
%   \item A mandatory Vulkan route is fail-closed. An independent GPU prefilter
%   remains a local acceleration inside a CPU or SIMD pipeline.
% \end{enumerate}

% Several limits remain. The CPU stage chain and the resident GPU batch route
% share algorithms and result contracts, but they are not the same instruction
% graph. Direct use of a persisted GPU index currently supports coordinate
% projection, not persisted CIGAR extension. Final result objects are built on
% the host. Small query groups may also be slower on a GPU because submission,
% allocation, and transfer costs are not amortized.

% Within these limits, BioSeqDB turns a procedural mapper into a bounded database
% operator. It retains the seeding, chaining, and extension structure
% of minimap2 while moving target preparation outside repeated query work.
% Typed stages define the data exchanged by algorithms. Lazy sequence access,
% memory-admitted batching, controlled CPU parallelism, and resident Vulkan
% waves connect those algorithms to PostgreSQL execution. The following
% chapters examine the principal biological stages in detail, beginning with
% minimizer extraction.
